\section{Multimodal Hallucination}
\label{sec:hallu}
Multimodal hallucination refers to the phenomenon of responses generated by MLLMs being inconsistent with the image content~\cite{yin2023woodpecker}. As a fundamental and important problem, the issue has received increased attention. In this section, we briefly introduce some related concepts and research development.

\subsection{Preliminaries}
\label{sec:hallu_prelim}
Current research on multimodal hallucinations can be further categorized into three types~\cite{zhai2023halle}:
\begin{enumerate}
    \item \textit{Existence Hallucination} is the most basic form, meaning that models incorrectly claim the existence of certain objects in the image.
    \item \textit{Attribute Hallucination} means describing the attributes of certain objects in a wrong way, \eg failure to identify a dog's color correctly. It is typically associated with existence hallucination since descriptions of the attributes should be grounded in objects present in the image.
    \item \textit{Relationship Hallucination} is a more complex type and is also based on the existence of objects. It refers to false descriptions of relationships between objects, such as relative positions and interactions.   
\end{enumerate}

In what follows, we first introduce some specific evaluation methods (\S \ref{sec:hallu_eval_method}), which are useful to gauge the performance of methods for mitigating hallucinations (\S \ref{sec:hallu_method}). Then, we will discuss in detail the current methods for reducing hallucinations, according to the main categories each method falls into.

\subsection{Evaluation Methods}
\label{sec:hallu_eval_method}

CHAIR~\cite{rohrbach2018object} is an early metric that evaluates hallucination levels in open-ended captions. The metric measures the proportion of sentences with hallucinated objects or hallucinated objects in all the objects mentioned.
In contrast, POPE~\cite{li2023evaluating} is a method that evaluates closed-set choices. Specifically, multiple prompts with binary choices are formulated, each querying if a specific object exists in the image. The method also covers more challenging settings to evaluate the robustness of MLLMs, with data statistics taken into consideration. The final evaluation uses a simple watchword mechanism, \ie by detecting keywords ``yes/no'', to convert open-ended responses into closed-set binary choices.
With a similar evaluation approach, MME~\cite{mme} provides a more comprehensive evaluation, covering aspects of existence, count, position and color, as exemplified in~\cite{yin2023woodpecker}.

Different from previous approaches that use matching mechanisms to detect and decide hallucinations, HaELM~\cite{wang2023evaluation} proposes using text-only LLMs as a judge to automatically decide whether MLLMs' captions are correct against reference captions. In light of the fact that text-only LLMs can only access limited image context and require reference annotations, Woodpecker~\cite{yin2023woodpecker} uses GPT-4V to directly assess model responses grounded in the image.
FaithScore~\cite{jing2023faithscore} is a more fine-grained metric based on a routine that breaks down descriptive sub-sentences and evaluates each sub-sentence separately.
Based on previous studies, AMBER~\cite{wang2023llm} is an LLM-free benchmark that encompasses both discriminative tasks and generative tasks and involves three types of possible hallucinations (see \S \ref{sec:hallu_prelim}).

\subsection{Mitigation Methods}
\label{sec:hallu_method}
According to high-level ideas, the current methods can be roughly divided into three categories: pre-correction, in-process-correction, and post-correction.

\noindent \textbf{Pre-correction.}
An intuitive and straightforward solution for hallucination is to collect specialized data (\eg negative data) and use the data for fine-tuning, thus resulting in models with fewer hallucinated responses. 

LRV-Instruction~\cite{liu2024mitigating} introduces a visual instruction tuning dataset. Apart from common positive instructions, the dataset incorporates delicately designed negative instructions at different semantic levels to encourage responses faithful to the image content.
LLaVA-RLHF~\cite{sun2023aligning} collects human-preference pairs and finetunes models with reinforcement learning techniques, leading to models more aligned with less hallucinated answers.

\noindent \textbf{In-process-correction.}
Another line is to make improvements in architectural design or feature representation. These works try to explore the reasons for hallucinations and design corresponding remedies to mitigate them in the generation process.

HallE-Switch~\cite{zhai2023halle} performs an empirical analysis of possible factors of object existence hallucinations and hypothesizes that existence hallucinations derive from objects not grounded by visual encoders, and they are actually inferred based on knowledge embedded in the LLM. Based on the assumption, a continuous controlling factor and corresponding training scheme are introduced to control the extent of imagination in model output during inference. 

VCD~\cite{leng2023mitigating} suggests that object hallucinations derive from two primary causes, \ie statistical bias in training corpus and strong language prior embedded in LLMs. The authors take notice of the phenomenon that when injecting noise into the image, MLLMs tend to lean towards language prior rather than the image content for response generation, leading to hallucinations. Correspondingly, this work designs an amplify-then-contrast decoding scheme to offset the false bias.

HACL~\cite{jiang2023hallucination} investigates the embedding space of vision and language. Based on the observation, a contrastive learning scheme is devised to pull paired cross-modal representation closer while pushing away non-hallucinated and hallucinated text representation.

\noindent \textbf{Post-correction.}
Different from previous paradigms, post-correction mitigates hallucinations in a post-remedy way and corrects hallucinations after output generation. 
Woodpecker~\cite{yin2023woodpecker} is a training-free general framework for hallucination correction. Specifically, the method incorporates expert models to supplement contextual information of the image and crafts a pipeline to correct hallucinations step by step. The method is interpretable in that intermediate results of each step can be checked, and objects are grounded in the image. 
The other method LURE~\cite{zhou2023analyzing} trains a specialized revisor to mask objects with high uncertainty in the descriptions and regenerates the responses again.
